\section{The integral branching map and its exact normalization}
\label{ib-section}

Multiplicity one does not determine an integral normalization.
We prove integrality for the exact scalar used by the submitted
lowering map.  The ingredients are the preceding branching law,
the integral spanning theorem for initial minors
(Theorem~\ref{sm-initial-minors}), and the Capelli identities
of Section~\ref{rp-section}.  The central differential operators
will act on rows; this gives a direct projection argument.

First let $N=n+1$, $n\ge1$, let
$\lambda=(\lambda_1,\ldots,\lambda_n,0)$ be a partition,
and suppose $\lambda_i\ge\mu_i\ge\lambda_{i+1}$.
Use the polynomial modules and normalized Fock pairings
of Section~\ref{coef:section}.  Set
\[
 a_i=\mu_i-\lambda_{i+1},\quad b_i=\lambda_i-\mu_i,\quad
 d=\sum_i b_i,\quad D_0=\sum_i a_i,\quad
 D_1=\sum_i\min(\lambda_i-\lambda_{i+1},D_0).
 \tag{IB1}
\]
Define integer intervals and an integer $c$ by
\[
 \begin{split}
 R_i&=[\lambda_{i+1},\min(\lambda_i,\lambda_{i+1}+D_0)]\cap\mathbb Z,\\
 c&=\prod_{i=1}^n\ \prod_{v\in R_i\setminus\{\mu_i\}}
                  \prod_{j=1}^n\bigl((v-i)-(\mu_j-j)\bigr).
 \end{split}
 \tag{IB2}
\]
Empty products are one.  This is exactly $c_{\rm Br}$ in the
source: changing indices from zero to one cancels in each
difference.  Put $\xi=\xi_{\lambda,\mu}$ from (CB4), and let
$J:V_\mu\to V_\lambda$ be the intertwiner with $JH_\mu=\xi$.

\subsection{An integral lift of the required coefficient}

Write $\omega(\beta)=\sum_i i\beta_i$ for the weight of a
last-row monomial $z^\beta$, and write $[z^b]G$ for its
coefficient, a polynomial in the top block $A$.

\begin{lemma}\label{ib-integral-lift}
For every $f\in\nabla_\mu(\mathbb Z)$ there exists
$G\in\nabla_\lambda(\mathbb Z)$ of last-row degree $d$ with
\[
 G(A,z)=z^b f(A)+G_{<}(A,z),\qquad
 \omega(\beta)<\omega(b)
       \quad\text{for every last-row monomial }z^\beta\text{ of }G_{<}.
 \tag{IB3}
\]
In particular $[z^b]G=f$.
\end{lemma}
\begin{proof}
By Theorem~\ref{sm-initial-minors}, it suffices to lift a
product of initial-column minors of shape $\mu$; integer
linear combinations then give the assertion.
Induct on the number of nonconstant factors.
If there are none, $\mu=0$ and interlacing gives
$\lambda_2=\cdots=\lambda_N=0$.
The lift of $1$ is $z_1^{\lambda_1}$.

Choose a factor $M_I$ of height $s$, with rows
$I\subset\{1,\ldots,n\}$ and columns $1,\ldots,s$.
Thus $1\le s\le n$ and $\mu_s>\mu_{s+1}$, where
$\mu_{n+1}=0$.  Remove this factor to get shape
$\mu'=\mu-(1^s,0^{n-s})$.
If $\mu_s>\lambda_{s+1}$, subtract the same $s$
initial ones from $\lambda$ to obtain $\lambda'$.
Then $\lambda'$ interlaces $\mu'$: the only changed
lower boundary is $\lambda_{s+1}\le\mu_s-1$.
All subtracted entries are positive.  The other
inequalities follow by subtracting one from both
sides or leaving both sides unchanged.
Here $b'=\lambda'-\mu'=b$.
Multiply the inductively obtained lift by $M_I(A)$.

If $\mu_s=\lambda_{s+1}$, then $s<n$:
a height-$n$ factor would give $\mu_n\ge1>\lambda_N$.
Subtract $s+1$ initial ones from $\lambda$ to obtain
$\lambda'$.  Its upper inequality at position $s+1$ is
\[
              \lambda_{s+1}-1=\mu_s-1\ge\mu_{s+1}.
\]
At position $s$ the lower inequality is equality;
all other inequalities follow as before.
Thus $\lambda'$ interlaces $\mu'$ and
$b'=b-e_{s+1}\ge0$.
Append row $N$ to $I$ and take the initial-column
minor of height $s+1$.  Its last-row Laplace expansion is
\[
 \widetilde M_I(A,z)
   =z_{s+1}M_I(A)+
       \sum_{j=1}^{s}(-1)^{s+1+j}z_jM_{I,j}(A),
 \tag{IB4}
\]
where $M_{I,j}$ uses columns $1,\ldots,s+1$ with
column $j$ removed.  The leading last-row term has
weight $s+1$, and every other term has smaller weight.
Multiply this minor by the inductive lift.
Last-row weights add, proving (IB3).

Every factor used is an initial-column minor on $N$
rows, with integral coefficients and the required
column invariance.  Its shape adds to $\lambda$ and
its last-row degree adds to $d$.  This also verifies
membership and degree in both cases.
\end{proof}

\subsection{The exact integral projection}

\begin{lemma}\label{ib-capelli-projection}
The integer $c$ is nonzero.  On the last-row degree-$d$
subspace of $V_\lambda$, an operator with integer
differential coefficients preserves
$\nabla_\lambda(\mathbb Z)$ and equals $c$ times the
projection onto the $V_\mu$ summand in (CB7).
\end{lemma}
\begin{proof}
Use the top-row operators $E_{ab}$ of (CB1), with
$1\le a,b\le n$.  They preserve integral coefficients,
column degrees, column invariance, and last-row degree.
With indices starting at one write
\[
 C(u)=\sum_{\sigma\in\mathfrak S_n}\operatorname{sgn}(\sigma)
       \prod_{j=1}^n\bigl((u+j)\delta_{\sigma(j),j}
                                         -E_{\sigma(j),j}\bigr).
 \tag{IB5}
\]
The product is ordered by increasing $j$.
This is $B_n(u+1)$ from Corollary~\ref{rp-negative-convention}.
It is central and acts on $V_\eta$ by
$\chi_\eta(u)=\prod_j(u+j-\eta_j)$.
Put
\[
 Z=\prod_{i=1}^n\ \prod_{v\in R_i\setminus\{\mu_i\}} C(v-i).
 \tag{IB6}
\]
Its factors commute and preserve the integral polynomial module.

An interlacing summand of last-row degree $d$ has
$|\eta|=|\mu|$, so
$\sum_i(\eta_i-\lambda_{i+1})=D_0$.
Each summand is nonnegative, hence every $\eta_i\in R_i$.
If $\eta\ne\mu$, choose $i$ with $\eta_i\ne\mu_i$.
The included factor $C(\eta_i-i)$ has scalar zero on
$V_\eta$.  On $V_\mu$ the product has scalar exactly $c$.

Fix an included $v\in R_i\setminus\{\mu_i\}$.
The $j=i$ factor in (IB2) is nonzero.
If $j<i$, then
$\mu_j\ge\lambda_{j+1}\ge\lambda_i\ge v$, and
$v-i<\mu_j-j$.
If $j>i$, then
$v\ge\lambda_{i+1}\ge\lambda_j\ge\mu_j$, and
$v-i>\mu_j-j$.
Every factor is nonzero.
The branching law now proves the projection assertion.
\end{proof}

\begin{theorem}\label{ib-iota-integral}
The map $I=cJ$ is equivariant, sends $H_\mu$ to $c\xi$,
has last-row degree $d$, and satisfies
\[
                    I\nabla_\mu(\mathbb Z)\subset\nabla_\lambda(\mathbb Z).
 \tag{IB7}
\]
\end{theorem}
\begin{proof}
Take the integral lift $G$ of $f$ from
Lemma~\ref{ib-integral-lift}, and let $P_\mu$ be the
projection in Lemma~\ref{ib-capelli-projection}.
Coefficient extraction commutes with every top-row operator.
On $f\in V_\mu$, the operator (IB6) acts as $c$, so
\[
                    [z^b]P_\mu G=c^{-1}[z^b]ZG=f.
\]
The restriction of $[z^b]$ to $J(V_\mu)$ is the inverse
of $J$: it is equivariant and sends $\xi$ to $H_\mu$,
by (CB4) and its Laplace expansion.
The uniqueness in Lemma~\ref{cb-highest-classification}
therefore gives $[z^b]J=1$.
Consequently $P_\mu G=Jf$ and $If=ZG$, which is integral.
This also proves independence of the chosen lift.
The remaining properties follow from those of $J$.
\end{proof}

\subsection{Adjoints, the standard formula, and the norm}

Let $\Pi_\lambda,\Pi_\mu$ be the orthogonal Fock
projections onto the polynomial modules.  They have
rational matrices: choose rational bases and invert
their positive definite rational Gram matrices.
Set
\[
 C_\lambda=[H_\lambda,H_\lambda],\quad
 C_\mu=[H_\mu,H_\mu],\quad
 \kappa=\frac{[\xi,\xi]}{C_\lambda}>0,\qquad b!=\prod_i b_i!.
 \tag{IB8}
\]

\begin{theorem}\label{ib-standard-integral}
On the indicated polynomial modules define
\[
 \begin{split}
 \mathcal B F&=c\kappa\,\Pi_\mu([z^b]F),\\
 \mathcal I f&=\frac{c[\xi,\xi]}{b!C_\mu}\Pi_\lambda(z^b f).
 \end{split}
 \tag{IB9}
\]
Then $\mathcal I=I=cJ$ and $\mathcal B$ is its adjoint
for the normalized Fock pairings.  In particular
\[
 \mathcal I\nabla_\mu(\mathbb Z)\subset\nabla_\lambda(\mathbb Z),
 \qquad
 \mathcal B\Delta_\lambda(\mathbb Z)\subset\Delta_\mu(\mathbb Z).
 \tag{IB10}
\]
Both maps are nonzero and equivariant.  Their squared
operator norms for the normalized Hermitian Fock
metrics are exactly $c^2\kappa$.
\end{theorem}
\begin{proof}
Coefficient extraction commutes with $\operatorname{GL}_n$.
The projection $\Pi_\mu$ commutes with $U(n)$ by
unitarity, hence with its complexified Lie algebra
and with $\operatorname{GL}_n(\mathbb C)$.
Thus $\phi=\Pi_\mu[z^b]$ vanishes on every last-row
degree other than $d$; in degree $d$, multiplicity
one makes it vanish on every summand except
$J(V_\mu)$, on which $\phi J=1$.

The Gram-matrix calculation in
Lemma~\ref{cb-highest-classification} shows that the
pullback of the normalized Hermitian form under $J$
is a scalar multiple of the form on $V_\mu$.
Evaluation on $H_\mu$ gives the scalar $\kappa$.
Orthogonality of the summands therefore gives
$J^*=\kappa\phi$ and $I^*=\mathcal B$.
Their exact squared norms are $c^2\kappa>0$.

Comparing monomials in the top and bottom rows gives
\[
                         [F,z^b f]=b!\,[[z^b]F,f].
\]
After normalizing, the adjoint of $\phi$ is
$C_\lambda(b!C_\mu)^{-1}\Pi_\lambda(z^b\,\cdot)$.
Taking the adjoint of $c\kappa\phi$ proves the second
formula of (IB9), including every factorial.
All matrices are rational, so the same adjoint
identity holds for the rational bilinear forms.

The first containment in (IB10) is (IB7).
For $F\in\Delta_\lambda(\mathbb Z)$ and
$f\in\nabla_\mu(\mathbb Z)$ it follows that
\[
 B_\mu(\mathcal B F,f)=B_\lambda(F,\mathcal I f)\in\mathbb Z.
\]
Theorem~\ref{coef:integral-duality} gives
$\mathcal B F\in\Delta_\mu(\mathbb Z)$.
No prime dividing $b!$ or $C_\mu$ has been excluded.
\end{proof}

\begin{proposition}\label{ib-norm-bound}
Write $M=\lambda_1-\lambda_N$.  Then
\[
 |c|\le(M+n-1)^{nD_1},\qquad
 \kappa\le\dim V_\lambda,\qquad
 \|\mathcal B\|=\|\mathcal I\|
                  \le |c|\sqrt{\dim V_\lambda}.
 \tag{IB11}
\]
A zero base in the first power has exponent zero
and value $1$.  In particular these estimates imply
the source bound
\[
 \|\mathcal B\|,\|\mathcal I\|
       \le\bigl(2(M+N)\bigr)^{ND_1}\sqrt{\dim V_\lambda}.
 \tag{IB12}
\]
\end{proposition}
\begin{proof}
The interval $R_i$ contains $\mu_i$, and
$\#(R_i\setminus\{\mu_i\})
=\min(\lambda_i-\lambda_{i+1},D_0)$.
Both $v$ and $\mu_j$ in (IB2) lie between $\lambda_N$
and $\lambda_1$.  Every factor has absolute value
at most $M+n-1$, and there are $nD_1$ factors.
If $M+n-1=0$, then $n=1,M=0$ and $D_1=0$.

Normalize Haar measure on $U(N)$ to have mass one.
The highest vector $h=H_\lambda$ has normalized norm one.
Average the rank-one projection onto $\rho(g)h$.
The resulting self-adjoint operator commutes with
the group, so each of its eigenspaces is invariant.
Irreducibility makes it a scalar operator.
Its trace is one, hence its scalar is $(\dim V_\lambda)^{-1}$.
Consequently
\[
 \int_{U(N)}|\langle\xi,\rho(g)h\rangle_\lambda|^2\,dg
                         =\frac{\kappa}{\dim V_\lambda}.
\]
The evaluation identity of
Lemma~\ref{coef:extremal-weight} and unitarity identify
this matrix coefficient, up to conjugation and
inversion of $g$, with $\xi(g^t I_{N,n})$.
Each initial or replaced-row minor of a matrix with
orthonormal columns has absolute value at most one:
it is a coordinate of the unit exterior product of
those initial columns.  Formula (CB4) therefore
gives $|\xi(g^tI_{N,n})|\le1$.
The integral is at most one, proving the bound
on $\kappa$.  The exact norm formula proves the
last part of (IB11).  Finally $M+N\ge2$ and $n\le N$
give (IB12).
\end{proof}

\subsection{The local basis identity and a necessary scalar}

The statements extend to coefficient fields and tensor
products over embeddings.  The integral containments
also extend to characteristic-zero coefficient domains:
initial minors span the dual-Weyl module, the lift
and $Z$ have integer coefficients, and the perfect
Weyl duality extends with the coefficient ring.
In particular (IB10) imposes no exclusion of small primes.

For dominant integral $\lambda,\mu$ with the same
smallest entry $s$, subtract $s$ from both tuples.
Their canonical polynomial modules are those above;
the maps and integral pairings do not change.
The action acquires the common factor $\det^s$,
and the last-circle character becomes
$t^{d+s}=t^{|\lambda|-|\mu|}$.

\begin{corollary}\label{ib-diagonal-covariance}
Let $\lambda,\mu$ interlace and have the same smallest
entry.  Put $d_0=|\lambda|-|\mu|$ and
$\mathcal B_t=t^{-d_0}\mathcal B$ for $t\ne0$.
For $R_\lambda(g)=\rho_\lambda(g^{-t})$ and
$B=\operatorname{diag}(B',t^{-1})$ one has
\[
                R_\mu(B'^{-1})\mathcal B_t
                              =\mathcal B R_\lambda(B^{-1}).
 \tag{IB13}
\]
For tensor products the scalar is the product of
the corresponding scalars over the embeddings.
\end{corollary}
\begin{proof}
Equivariance and the last-circle character give
$\mathcal B\rho_\lambda(\operatorname{diag}(g,a))
=a^{d_0}\rho_\mu(g)\mathcal B$.
Now $R_\lambda(B^{-1})=\rho_\lambda(B^t)$ and
$B^t=\operatorname{diag}(B'^t,t^{-1})$.
Substitution proves the formula; tensor products
prove the final assertion.
In the arithmetic application $t=\overline{\pi_T}$,
as dictated by the Hermitian dual lattice.
\end{proof}

The scalar $c$ cannot simply be deleted.
Take $\lambda=(2,1,0)$ and $\mu=(2,0)$.
Write the two columns as $x=(x_1,x_2,x_3)^t$ and
$y=(y_1,y_2,y_3)^t$, and put $D_{ij}=x_i y_j-x_j y_i$.
Then $\xi=x_1D_{13}$ and
\[
 E_{21}\xi=x_1D_{23}+x_2D_{13}
                         =2x_2D_{13}-x_3D_{12}.
\]
Since $E_{21}(x_1^2)=2x_1x_2$, the unscaled map gives
\[
                  J(x_1x_2)=x_2D_{13}-\tfrac12x_3D_{12}.
\]
The coefficient of $x_1x_3y_2$ is $-1/2$.
Thus this is not an integral polynomial, although
$x_1x_2$ is in the integral dual-Weyl module.
Here $D_0=1$, $R_1=\{1,2\}$, $R_2=\{0,1\}$, and
(IB2) gives $c=(-2)(-2)=4$.
The actual map is $4J$, whose integrality is proved
in Theorem~\ref{ib-standard-integral}.
This refutes deletion of the normalization, not
the submitted integral branching assertion.

These results settle the polynomial branching input
and its local basis identity.  The actual twisted
automorphic classes, their coefficient lattices,
and their period sums are separate arithmetic statements.



% A mathematical section, to be included after the common preamble.
% Requires amsmath, amssymb, amsthm, booktabs and the theorem environments
% theorem, proposition, lemma, corollary, example, remark.
% All labels in this file have the prefix div:.
